\documentclass[11pt,a4paper]{ctexart}
\usepackage[margin=2.25cm]{geometry}
\usepackage{amsmath,amssymb,booktabs,graphicx,hyperref,listings,microtype,siunitx,xcolor}
\usepackage{subcaption}
\usepackage{longtable}
\usepackage{fancyhdr}
\hypersetup{colorlinks=true,linkcolor=blue!55!black,urlcolor=blue!55!black,citecolor=blue!55!black}
\graphicspath{{results/figures/}}
\pagestyle{fancy}
\fancyhead[L]{CS336 Assignment 2: Systems}
\fancyhead[R]{ShaneLiu04}
\fancyfoot[C]{\thepage}
\setlength{\headheight}{14pt}
\setlength{\emergencystretch}{3em}
\sisetup{detect-all}
\lstset{basicstyle=\ttfamily\small,breaklines=true,frame=single}
\newcommand{\measured}{\textbf{[实测]}}
\newcommand{\derived}{\textbf{[理论推导]}}
\newcommand{\limited}{\textbf{[资源受限]}}

\newcommand{\missingfigure}[2]{%
  \begin{figure}[htbp]\centering
  \fbox{\parbox[c][3.2cm][c]{0.82\linewidth}{\centering
  #2\\[0.5em]\small 预期产物：\texttt{\detokenize{#1}}\\
  报告资产缺失；请运行 \texttt{scripts/build\_report\_assets.py}。}}
  \caption{#2（生成资产缺失）}
  \end{figure}}
\newcommand{\resultfigure}[3]{%
  \IfFileExists{results/figures/#1}{\begin{figure}[htbp]\centering
  \includegraphics[width=0.96\linewidth]{results/figures/#1}
  \caption{#2}\label{#3}\end{figure}}{\missingfigure{#1}{#2}}}

\title{\textbf{CS336 Assignment 2: Systems}\\
\large 从算子级优化到分布式训练的可复现实验}
\author{\href{https://github.com/ShaneLiu04}{ShaneLiu04}}
\date{\today}

\begin{document}
\maketitle

\begin{abstract}
本文实现并分析 Transformer 系统优化中的四条主线：基于 online softmax 的
FlashAttention、通信与反向计算重叠的 Distributed Data Parallel（DDP）、
Optimizer State Sharding，以及教学型 Fully Sharded Data Parallel（FSDP）。
所有数值结果均由版本化脚本写入原始 CSV，再自动生成图表；无法在单张 RTX
6000D 上完成的多卡实验被明确标注，而不以模拟结果代替实测。实测中 xl 模型的 BF16
forward+backward 为 438.12 ms，相比 FP32 的 1120.75 ms 加速 2.56 倍；
Triton forward 在 $N=8192,d=64$ 的 BF16 配置上为 0.307 ms，相比 PyTorch
SDPA 的 3.122 ms 加速 10.17 倍，同时把 peak allocated memory 从
853.13 MiB 降至 12.16 MiB；新增 fused backward 后，同配置完整
forward+backward 为 0.796 ms，相比 SDPA 加速 6.79 倍。报告由 228 条性能记录、
256 条数值误差记录、12 条 checkpointing 记录和 14 组解释图片支撑。
\end{abstract}

\tableofcontents
\newpage

\section{目标、范围与可复现性}
\subsection{实验原则}
计时前执行 warmup，并在 CUDA 计时边界调用同步；报告均值与标准差，而不是最优单次值。
OOM 同样作为结果保留。每条记录包含 commit SHA、GPU、PyTorch/CUDA 版本、dtype、
batch、sequence length、warmup 和 repetitions。入口为
\texttt{scripts/benchmark\_systems.py}，绘图入口为
\texttt{scripts/build\_report\_assets.py}。与 Assignment 1 相同，所有报告资产集中于
\texttt{report/results/}：\texttt{raw/} 保存不可变原始指标，\texttt{figures/} 同时保存
PNG/PDF，\texttt{tables/} 保存可复用 LaTeX 表格，\texttt{summary.json} 汇总关键结论，
\texttt{manifest.sha256} 记录逐文件校验值。

\subsection{硬件边界}
可用远端为单张 RTX 6000D（85,651 MiB，driver 595.71.05）。实测环境为
PyTorch 2.8.0+cu128 与 Triton 3.4.0；由于锁文件中的 2.11.0 wheel 在实验窗口内下载
持续停滞，本文没有把 2.8 的结果标成 2.11。该硬件足以验证 Triton kernel、单卡显存与端到端性能，但不能
产生课程指定的 2/4/6-GPU NCCL 曲线、2-GPU DDP/FSDP overlap trace 或 2$\times$B200
Leaderboard 成绩。正确性测试中的两个 rank 使用 Gloo 多进程；这验证数学语义，但不代表
NCCL 的真实性能。本文始终区分``实测''、``正确性验证''和``理论估算''。

\section{基准模型与测量方法}
除算子 microbenchmark 外，基准保持 handout 的 vocabulary size 10,000、batch size 4
和默认 context length 512。模型规模见表~\ref{tab:model-config}。
\begin{table}[htbp]\centering
\caption{Transformer 配置}\label{tab:model-config}
\begin{tabular}{lrrrr}
\toprule
规模 & $d_{\mathrm{model}}$ & $d_{\mathrm{ff}}$ & 层数 & heads\\
\midrule
small & 768 & 3072 & 12 & 12\\
medium & 1024 & 4096 & 24 & 16\\
large & 1280 & 5120 & 36 & 20\\
xl & 2560 & 10240 & 32 & 32\\
10B & 4608 & 12288 & 50 & 36\\
\bottomrule
\end{tabular}
\end{table}

不进行 warmup 会同时混入 lazy library loading、CUDA context 建立、allocator 扩容和
kernel autotuning，因而明显高估稳定迭代成本。对 compiled graph，还会混入 compilation
成本。由于本次需要在有限租用窗口内同时覆盖 182 个 performance cases，统一使用
2 次 warmup、5 次测量；原始 CSV 同时保存 standard deviation，避免把单次最优值当结论。

\section{Profiling、混合精度与显存}
\subsection{时间分解}
一个训练 step 可分为 forward、loss、backward、gradient synchronization 和 optimizer
step。NVTX range 应包住测量区间，并把 warmup 排除在 trace 外。随着 sequence length
增大，标准 attention 的 $QK^\top$ 与 $PV$ 逐渐主导；模型变宽后 GEMM 占比上升，
而小模型更容易受 launch overhead 和 Python overhead 影响。

\subsection{Mixed precision}
BF16 保留与 FP32 相同的 8-bit exponent，因此相较 FP16 更不易 overflow；mantissa
缩短会引入舍入误差。归一化和 reduction 的累加精度尤其重要：若把大量小量直接累加到
BF16，低位信息会持续丢失。因此实现保留 FP32 master weights，并在 softmax 统计量、
loss 和必要 reduction 上使用 FP32。

\resultfigure{transformer_scaling.pdf}
{RTX 6000D 上模型规模、精度与 step latency/peak memory 的关系。每点为 2 次 warmup
后 5 次测量的均值；误差来源见原始 CSV。}{fig:transformer-scaling}

本次实际 sweep 使用 2 次 warmup 与 5 次测量，以在有限租用窗口内覆盖全部四种模型与
两种精度；图~\ref{fig:transformer-scaling} 的原始记录保留了该元数据。small/medium/large/xl
的 BF16 step time 分别为 35.87/94.06/190.51/438.12 ms，FP32 则为
58.62/177.26/411.86/1120.75 ms。BF16 优势随模型规模增大而扩大；xl 的 peak allocated
memory 从 40,990 MiB 降至 38,058 MiB，说明该 workload 中 FP32 master-like 参数与其他
常驻项占比较高，autocast 并不会把总显存简单减半。

\resultfigure{mixed_precision_benefit.pdf}
{BF16 相对 FP32 的实测 speedup 与 peak-memory reduction。模型越大，Tensor Core
吞吐收益越明显；但 xl 的显存降幅仅 7.2\%，说明常驻 FP32 状态占比上升。}
{fig:mixed-precision}

\measured 图~\ref{fig:mixed-precision} 给出两条不同趋势：speedup 从 small 的 1.63 倍
单调上升到 xl 的 2.56 倍，而显存降幅从 22.2\% 下降到 7.2\%。这并不矛盾：BF16 让大型
GEMM 更接近 Tensor Core 的高吞吐区间，但本 benchmark 未引入 BF16 optimizer states，
且参数、gradient、allocator reserve 等常驻部分不会按 activation dtype 等比例缩小。
因此 mixed precision 的首要收益是算力利用率，而不是承诺``显存减半''。

\subsection{Context scaling 与 \texttt{torch.compile}}
\resultfigure{model_context_scaling.pdf}
{Small/medium 模型在 BF16 下从 sequence 128 扩展到 2048 的 forward+backward latency
和 peak memory。}{fig:context-scaling}
\resultfigure{torch_compile_model.pdf}
{Sequence 512 上 eager 与 \texttt{torch.compile(mode=reduce-overhead)} 的时间和显存。}
{fig:torch-compile}

\measured Medium 模型从 sequence 512 增至 2048 时，step latency 从 90.76 ms 增至
998.89 ms（11.0 倍），peak memory 从 8.36 GiB 增至 56.68 GiB（6.78 倍）。
增长明显快于 sequence 的 4 倍，反映标准 attention 的 quadratic component 逐渐主导。
Small 模型同区间为 39.92$\rightarrow$374.10 ms，显存
3.18$\rightarrow$22.77 GiB。

\measured Compile 对 small FP32/BF16 的 steady-state speedup 为 1.20/1.98 倍，
对 medium 为 1.25/1.72 倍。BF16 medium 的 peak memory 从 8.38 GiB 降至
7.14 GiB。该结果排除了首次 compilation latency（2 次 warmup），因此只代表重复 step；
短任务若无法摊销 graph compilation，wall-clock 未必获益。

\subsection{显存核算}
训练显存由参数、梯度、optimizer state、activation 和临时 workspace 构成。以 AdamW
的 FP32 训练为例，忽略 allocator fragmentation 时，每参数约包含
\[
4\text{ B (weight)}+4\text{ B (gradient)}+
8\text{ B (first/second moments)}=16\text{ B}.
\]
BF16 autocast 主要压缩 activation 和算子输入；若 master weights、gradients 和 Adam
states 仍为 FP32，参数相关部分不会简单减半。memory snapshot 中的峰值还需结合 allocation
stack 判断，不能只凭张量理论大小解释。

\section{Activation Checkpointing}
设网络有 $L$ 个等成本 block。保存每个 block 的边界 activation、在 backward 中重算
block 内部 forward，可把 block 内 activation 从随 $L$ 累积改为只保留边界，代价是约一次
额外 forward。更一般地，每 $k$ 层保存 checkpoint，存储量近似
$O(L/k+k)$，在 $k\approx\sqrt{L}$ 时达到 $O(\sqrt L)$。
实现中应 checkpoint 纯函数式 block，并显式处理 RNG state，避免 dropout 重算不一致。
由于 xl、batch 4、sequence 2048 会让比较受 OOM 边界支配，受控实验采用 large、
batch 2，并扫描 sequence 512/1024/2048 与 0/4/12/36 个 checkpoint segments。

\resultfigure{activation_checkpointing.pdf}
{RTX 6000D 上 large 模型的 activation checkpointing trade-off。左图是重算时间，
右图是 peak allocated memory；0 表示禁用 checkpoint。}{fig:checkpointing}

\measured 在 sequence 2048 上，禁用 checkpoint 时 step 为 983.81 ms、峰值
56.48 GiB；36 segments 时为 1254.51 ms、11.53 GiB。即用 27.5\% latency overhead
换取 79.6\% peak-memory reduction。最初的 4 segments 已把峰值降至 21.84 GiB，
继续增加到 12/36 segments 仍能降至 14.11/11.53 GiB，但时间基本平台化在
1.25 s 左右。这说明 checkpoint 粒度的选择存在收益递减：如果 21.84 GiB 已满足设备约束，
4 segments 的 22.4\% overhead 比 36 segments 更合理；只有在显存是硬约束时才应选择
最细粒度。

\section{Attention 与 FlashAttention-2}
\subsection{Online softmax}
标准 attention 为
\[
S=QK^\top/\sqrt d,\qquad P=\operatorname{softmax}(S),\qquad O=PV.
\]
它会物化 $N\times N$ 的 $S$ 或 $P$。FlashAttention 按 key/value tile 流式更新每行最大值
$m$、归一化因子 $\ell$ 和未归一化输出 $o$。新 tile 的统计量为
\[
m'=\max(m,m_j),\quad
\ell'=e^{m-m'}\ell+\sum_t e^{S_{jt}-m'},\quad
o'=e^{m-m'}o+\sum_t e^{S_{jt}-m'}V_t.
\]
最终输出 $O=o/\ell$，log-sum-exp 为 $L=m+\log\ell$。该重标定避免 overflow，
同时把额外显存从 $O(N^2)$ 降为线性量级。

\subsection{Backward}
令 $D_i=\sum_d O_{id}\,dO_{id}$。对每个 tile 重算
$P_{ij}=\exp(S_{ij}-L_i)$，则
\[
dV=P^\top dO,\quad dS=P\odot(dO\,V^\top-D),\quad
dQ=dS\,K/\sqrt d,\quad dK=dS^\top Q/\sqrt d.
\]
因而 backward 无需保存完整 $P$。causal 模式在分数进入指数运算前屏蔽未来位置，避免
无效位置参与 online statistics。

\subsection{Forward 性能与 IO 收益}
\resultfigure{attention_forward_fp32.pdf}
{FP32 attention forward sweep。Triton kernel、PyTorch SDPA 与 Python tiled
online-softmax 使用相同输入；虚实线分别为 $d=64/128$。}{fig:attn-forward-fp32}
\resultfigure{attention_forward_bf16.pdf}
{BF16 attention forward sweep。Triton 的线性额外存储与标准 SDPA 的二次中间量形成
明显对照。}{fig:attn-forward-bf16}

\measured 在 $N=8192,d=64$ 的 BF16 forward 中，Triton 为 0.307 ms，PyTorch SDPA
为 3.122 ms，获得 10.17 倍 speedup；peak allocated memory 分别为 12.16 MiB 与
853.13 MiB，相差 70.2 倍。FP32 下 Triton/SDPA 分别为 0.364/2.801 ms，仍有 7.69 倍
speedup，显存为 16.16/850.13 MiB。随着 $N$ 从 4096 增至 8192，SDPA 的显存约扩大
3.8 倍，而 Triton 仅从 10--12 MiB 增至 12--16 MiB，直接验证了避免物化 $N^2$
attention matrix 的 IO 优势。

Python tiled 版本在 $N=8192,d=64$ 上需要 2369.66 ms（BF16），比 Triton 慢约
7717 倍。这不是 online-softmax 算法本身的问题，而是 host 端双重 tile loop 触发大量
小 kernel launch；同一数学分块只有融合进一个 persistent Triton program 后才具备性能意义。

\subsection{优化前基线：未融合 backward}
\resultfigure{attention_end_to_end_fp32.pdf}
{FP32 attention forward+backward sweep。Triton 路径的 backward 当前为 tiled PyTorch
重计算，因此该图隔离了尚未融合的性能缺口。}{fig:attn-e2e-fp32}
\resultfigure{attention_end_to_end_bf16.pdf}
{BF16 attention forward+backward sweep。}{fig:attn-e2e-bf16}
\resultfigure{attention_latency_memory_tradeoff.pdf}
{$N=4096,d=64$ 下 latency--memory Pareto 对比。左/下更优，但当前没有实现同时占据
左下角。}{fig:attn-tradeoff}

\measured 在最初的 $N=4096,d=64$ FP32 forward+backward 中，SDPA、PyTorch tiled 与
Triton 路径分别为 1.10、1117.99 与 556.09 ms；对应 peak allocated memory 为
280.25、24.38 与 24.38 MiB。Triton 路径节省约 11.5 倍显存，但被 tiled PyTorch
backward 主导，因此该历史版本尚未形成 end-to-end speedup。图~\ref{fig:attn-tradeoff} 显示：
SDPA 位于``快但高显存''端，两个 tiled 路径位于``低显存但慢''端。下一优化优先级必须是
Triton backward fusion，而不是继续微调已经很快的 forward tile size。

\subsection{Fused Triton backward}
\resultfigure{fused_triton_end_to_end.pdf}
{实现 delta preprocess、$dQ$ 与 $dK/dV$ 三阶段 Triton kernels 后的完整
forward+backward latency 与 peak memory。}{fig:fused-e2e}
\resultfigure{triton_backward_speedup.pdf}
{相对旧 PyTorch tiled backward fallback 的 end-to-end speedup。}{fig:fused-speedup}

\measured 在 BF16、$N=4096,d=64$ 上，fused 路径从 597.05 ms 降至 0.320 ms，
加速 1864.7 倍，并比 PyTorch SDPA 的 1.144 ms 快 3.57 倍；peak memory 为
20.28 MiB，对比 SDPA 的 277.75 MiB。$N=8192$ 时 fused Triton 为 0.796 ms，
SDPA 为 5.404 ms，即 6.79 倍 speedup，同时把 peak memory 从 1051.25 MiB
降至 24.31 MiB。优化消除了 Python launch overhead，使 IO-aware forward 的优势真正
延伸到 end-to-end。
\measured 长上下文扩展到 $N=32768$ 时，fused Triton 为 10.06 ms、48.5 MiB，
而 SDPA 为 88.56 ms、16,444 MiB：速度提升 8.80 倍，peak allocated memory 相差
339 倍。该结果同时展示了 quadratic materialization 与 linear auxiliary state 在长序列下
逐渐扩大的差距。

\subsection{数值误差}
\resultfigure{attention_correctness.pdf}
{跨 causal/non-causal、sequence 128--1024、$d=64/128$、FP32/BF16 的 256 条
output/gradient 误差记录；纵轴为相对 PyTorch SDPA 的最大绝对误差。}
{fig:attn-correctness}

\measured PyTorch tiled FP32 的最大绝对误差为 $9.5\times10^{-6}$；Triton FP32
为 $9.0\times10^{-3}$，仍在官方测试的 $10^{-2}$ absolute tolerance 内。BF16 两条路径
的 envelope 均为 0.015625，等于该量级下 BF16 quantization step 的典型尺度。误差主要
出现在 $dK/dV$ reduction；对接近零的 reference 值，relative error 会被放大，因此报告
以 max absolute error 和官方 tolerance 为主要判断标准。

\section{Distributed Data Parallel}
\subsection{Naive、flat 与 overlap}
Naive DDP 在 backward 完成后逐参数 all-reduce，关键路径包含全部通信；flat-gradient
方案减少 collective launch 数，但仍可能把通信集中到 backward 末尾。本文实现为每个唯一
参数注册 post-accumulate hook：gradient ready 后立即发起 asynchronous all-reduce，
在 optimizer step 前统一等待 handle。共享参数按对象 identity 去重，避免 tied weights
被重复规约。

若每 rank 的 local loss 使用 mean reduction，则 all-reduce 后还需除以 world size：
\[
g=\frac1p\sum_{r=0}^{p-1}g_r.
\]
这与把所有 rank 的 batch 拼接后计算全局 mean loss 等价。初始化时从 rank 0 广播参数和
buffer，确保不同 seed 的进程从同一点开始。

\subsection{通信模型}
ring all-reduce 对 $M$ bytes、$p$ ranks 的每 rank 通信量约为
$2\frac{p-1}{p}M$，时间模型
\[
T_{\mathrm{AR}}\approx 2(p-1)\alpha+
2\frac{p-1}{p}\frac{M}{\beta},
\]
其中 $\alpha$ 是每步 latency，$\beta$ 是有效 bandwidth。小 tensor 受 latency 主导，
大 tensor 受 bandwidth 主导；这解释了 flat buffer 在小参数众多时可能获益。

\section{Optimizer State Sharding}
每个全局唯一参数按稳定索引分配 owner rank，只有 owner 构造底层 optimizer state 并执行
update；随后 owner 广播更新后的参数。这样 optimizer state 由每 rank 的 $O(P)$ 降为
约 $O(P/p)$，但参数和 gradient 仍复制，因此对应 ZeRO Stage 1，而不是 FSDP。

实现必须特别处理三点：(1) tied parameter 只出现一次；(2) \texttt{add\_param\_group}
在所有 rank 上保持相同 owner mapping；(3) checkpoint 要能恢复分片 state。广播位于
critical path，且按参数广播会受 collective latency 影响；生产实现通常 flatten/bucketize。

\section{Fully Sharded Data Parallel}
FSDP 进一步分片参数、gradient 与 optimizer state。forward 前 all-gather 当前 layer 的
full parameter，计算后释放；backward 中再次 materialize，并用 reduce-scatter 产生本 rank
的 gradient shard。理想状态下参数相关常驻显存从 $O(P)$ 降至 $O(P/p)$，代价是每层通信和
更严格的 lifetime 管理。

mixed precision 下通信/计算副本可用 FP16/BF16，但 optimizer 更新使用 FP32 master shard。
prefetch 下一层能隐藏一部分 all-gather；过早 prefetch 会增加同时驻留 full parameters
的数量，反而抬高峰值。正确的 Nsight 证据应显示 all-gather 位于对应 GEMM 前并尽量与前层
计算重叠。由于当前只有单卡，本报告不伪造该双卡 trace。

\section{并行策略推导}
\subsection{Data parallel 与 FSDP}
纯 data parallel 每 step 的主要通信是 gradient all-reduce，近似
$2(p-1)P/p$ 个元素。FSDP 对每层引入 parameter all-gather 与 gradient reduce-scatter；
总通信量与 DDP 同阶，但峰值常驻显存更低、collective 更细粒度。是否成为瓶颈取决于
\[
T_{\mathrm{comm,exposed}}>T_{\mathrm{compute,overlappable}}.
\]
因此只比较总 bytes 不足以判断性能，还必须考虑依赖关系、bucket 粒度与 fabric topology。

\subsection{Tensor parallel 与 2D parallel}
Tensor parallel 把矩阵乘的 contracting/non-contracting dimension 分到设备。column-parallel
linear 常在输出保持分片，row-parallel linear 则需要规约 partial outputs；反向传播的
collective 与前向布局互为对偶。把 TP 组大小记为 $t$、FSDP/DP 组大小记为 $f$，
总设备数 $p=tf$。TP 降低单层参数和 activation 的局部尺寸，但 collective 频率高；
FSDP 跨副本分片参数，collective 粒度较大。2D 组合应让高频 TP collective 留在高带宽域内，
让 FSDP 跨较慢链路，并以实测 topology 校正理论模型。

\section{测试、局限与后续工作}
自动测试覆盖 FlashAttention forward/backward、DDP 等价性与 tied weights、sharded
optimizer，以及 FSDP 的 FP32/FP16 correctness 和 gradient synchronization。测试应在
CPU/Gloo 与 CUDA/Triton 两条路径分别运行。单元测试通过只说明接口与小规模数学正确，
不自动证明 kernel 达到 IO-optimal，也不证明分布式实现具备生产级容错。

远端完成了 CUDA/Triton 6 项测试与 CPU/Gloo 分布式 8 项测试，合计 14/14 passed；
本地 PyTorch 2.11 CPU 环境另得到 10 passed、4 个 CUDA case skipped。当前最重要局限是
只有一张 GPU，无法提供 NCCL multi-GPU performance evidence。\path{scripts/run_remote_experiments.sh}
可重放测试、sweep 与绘图。若后续获得多卡环境，应补充 NCCL
all-reduce size sweep、naive/flat/overlap Nsight 对比、FSDP prefetch trace 和 2-GPU
accounting；在此之前这些项目保持``未实测''。

\resultfigure{test_matrix.pdf}
{本地 CPU、远端 CUDA/Triton 与远端 CPU/Gloo 三种环境的测试矩阵。远端两条路径分开
执行，避免单 GPU 上两个 Gloo rank 争用同一 CUDA device。}{fig:test-matrix}

\subsection{结论边界与可执行的下一步}
\begin{enumerate}
  \item \textbf{已完成的关键优化：}fused Triton backward 让 BF16、$N=4096$
  end-to-end 相比 fallback 加速 1864.7 倍，并开始同时优于 SDPA 的时间与显存。
  \item \textbf{Checkpoint 策略：}先按设备预算选择最粗可行粒度。large/sequence 2048
  上 4 segments 已节省 61.3\% 显存；36 segments 虽达到 79.6\%，但进一步增加重算且收益递减。
  \item \textbf{分布式证据：}当前 DDP/FSDP 只有 correctness evidence。获得至少两张 GPU
  后，应使用 NCCL + Nsight 量化 exposed communication；在此之前不报告伪造的 overlap speedup。
  \item \textbf{版本复验：}RTX 实验使用 PyTorch 2.8.0+cu128，而 assignment lock 指向
  2.11。升级后应重新生成 \texttt{summary.json} 与 \texttt{manifest.sha256}，而不是混用数字。
\end{enumerate}

\appendix
\section{报告资产与复现命令}
报告的每个定量结论都可追溯到 \texttt{report/results/raw/}。执行
\begin{lstlisting}
python scripts/benchmark_systems.py attention ...
python scripts/benchmark_correctness.py
python scripts/benchmark_checkpointing.py
python scripts/build_report_assets.py --results report/results
\end{lstlisting}
即可重建 12 组 PNG/PDF、3 个 LaTeX tables、\texttt{summary.json} 和
\texttt{manifest.sha256}。图像同时保留 PDF 是为了 LaTeX vector rendering，PNG
用于 README 预览；报告不手工抄写 plotting script 可以推导的表格。

\section*{致谢与使用声明}
本项目由 \href{https://github.com/ShaneLiu04}{ShaneLiu04} 完成，用于个人自学与公开学习记录；实验限制与
未完成项均如实披露。课程 handout 与 starter code 的版权归原作者所有。

\end{document}
